% GENERATED FILE. Edit canonical lessons, then run npm run generate:books.
% canonical-source-hash: fnv1a64:12fdcc614372d246
% canonical-lessons: PA-C146-beha, PA-C146-pakka, PA-C146-kacca, PA-C146-sacca, PA-C146-jhutha

\chapter{Stale, Ripe, Raw, True, False}
\label{ch:pa-stale-ripe-raw-true-false}
\hlchaptermodality{146}

\begin{chapteropening}
\textbf{By the end of this chapter:} \emph{I can say stale, ripe, raw, true and false.}

Complete the last lesson of chapter 146: I can say stale, ripe, raw, true and false.
\end{chapteropening}

\section[\texorpdfstring{\pa{ਬੇਹਾ} (behā)}{ਬੇਹਾ (behā)}]{\pa{ਬੇਹਾ} (behā) — stale}

\label{lesson:PA-C146-beha}



\begin{quote}
Before the new one: say the Punjabi for thirsty, then the Punjabi for fresh.
\end{quote}

\subsection*{You'll want to know: \pa{ਬੇਹਾ}}
\textbf{\pa{ਬੇਹਾ}} — \emph{behā} — \textquotedblleft{}stale\textquotedblright{}.

\begin{grammarlens}[title={What you've built}]
One more word for this chapter.
\end{grammarlens}

\subsection*{Your turn}
\begin{itemize}
\raggedright
  \item \emph{Say it:} \emph{behā}
  \item \emph{Say it:} \emph{behā}, once more
  \item \emph{Say it:} say \emph{behā}
  \item \emph{Recall:} say the Punjabi for thirsty, then the Punjabi for fresh, then say \emph{behā} again
\end{itemize}

\begin{tcolorbox}[breakable,colback=teal!4,colframe=teal!35!black,title={Before you move on}]
What does \pa{ਬੇਹਾ} mean? (\textquotedblleft{}Stale\textquotedblright{}.) Say it once more.
\end{tcolorbox}
\section[\texorpdfstring{\pa{ਪੱਕਾ} (pakkā)}{ਪੱਕਾ (pakkā)}]{\pa{ਪੱਕਾ} (pakkā) — ripe, firm}

\label{lesson:PA-C146-pakka}



\begin{quote}
Before the new one: say the Punjabi for fresh, then the Punjabi for stale.
\end{quote}

\subsection*{You'll want to know: \pa{ਪੱਕਾ}}
\textbf{\pa{ਪੱਕਾ}} — \emph{pakkā} — \textquotedblleft{}ripe, firm\textquotedblright{}.

\begin{grammarlens}[title={What you've built}]
One more word for this chapter.
\end{grammarlens}

\subsection*{Your turn}
\begin{itemize}
\raggedright
  \item \emph{Say it:} \emph{pakkā}
  \item \emph{Say it:} \emph{pakkā}, once more
  \item \emph{Say it:} say \emph{pakkā}
  \item \emph{Recall:} say the Punjabi for fresh, then the Punjabi for stale, then say \emph{pakkā} again
\end{itemize}

\begin{tcolorbox}[breakable,colback=teal!4,colframe=teal!35!black,title={Before you move on}]
What does \pa{ਪੱਕਾ} mean? (\textquotedblleft{}Ripe, firm\textquotedblright{}.) Say it once more.
\end{tcolorbox}
\section[\texorpdfstring{\pa{ਕੱਚਾ} (kaccā)}{ਕੱਚਾ (kaccā)}]{\pa{ਕੱਚਾ} (kaccā) — raw, unripe}

\label{lesson:PA-C146-kacca}



\begin{quote}
Before the new one: say the Punjabi for stale, then the Punjabi for ripe.
\end{quote}

\subsection*{You'll want to know: \pa{ਕੱਚਾ}}
\textbf{\pa{ਕੱਚਾ}} — \emph{kaccā} — \textquotedblleft{}raw, unripe\textquotedblright{}.

\begin{grammarlens}[title={What you've built}]
One more word for this chapter.
\end{grammarlens}

\subsection*{Your turn}
\begin{itemize}
\raggedright
  \item \emph{Say it:} \emph{kaccā}
  \item \emph{Say it:} \emph{kaccā}, once more
  \item \emph{Say it:} say \emph{kaccā}
  \item \emph{Recall:} say the Punjabi for stale, then the Punjabi for ripe, then say \emph{kaccā} again
\end{itemize}

\begin{tcolorbox}[breakable,colback=teal!4,colframe=teal!35!black,title={Before you move on}]
What does \pa{ਕੱਚਾ} mean? (\textquotedblleft{}Raw, unripe\textquotedblright{}.) Say it once more.
\end{tcolorbox}
\section[\texorpdfstring{\pa{ਸੱਚਾ} (saccā)}{ਸੱਚਾ (saccā)}]{\pa{ਸੱਚਾ} (saccā) — true}

\label{lesson:PA-C146-sacca}



\begin{quote}
Before the new one: say the Punjabi for ripe, then the Punjabi for raw.
\end{quote}

\subsection*{You'll want to know: \pa{ਸੱਚਾ}}
\textbf{\pa{ਸੱਚਾ}} — \emph{saccā} — \textquotedblleft{}true\textquotedblright{}.

\begin{grammarlens}[title={What you've built}]
One more word for this chapter.
\end{grammarlens}

\subsection*{Your turn}
\begin{itemize}
\raggedright
  \item \emph{Say it:} \emph{saccā}
  \item \emph{Say it:} \emph{saccā}, once more
  \item \emph{Say it:} say \emph{saccā}
  \item \emph{Recall:} say the Punjabi for ripe, then the Punjabi for raw, then say \emph{saccā} again
\end{itemize}

\begin{tcolorbox}[breakable,colback=teal!4,colframe=teal!35!black,title={Before you move on}]
What does \pa{ਸੱਚਾ} mean? (\textquotedblleft{}True\textquotedblright{}.) Say it once more.
\end{tcolorbox}
\section[\texorpdfstring{\pa{ਝੂਠਾ} (jhūṭhā)}{ਝੂਠਾ (jhūṭhā)}]{\pa{ਝੂਠਾ} (jhūṭhā) — false, lying}

\label{lesson:PA-C146-jhutha}



\begin{quote}
Before the new one: say the Punjabi for raw, then the Punjabi for true.
\end{quote}

\subsection*{You'll want to know: \pa{ਝੂਠਾ}}
\textbf{\pa{ਝੂਠਾ}} — \emph{jhūṭhā} — \textquotedblleft{}false, lying\textquotedblright{}.

\begin{grammarlens}[title={What you've built}]
One more word for this chapter.
\end{grammarlens}

\subsection*{Your turn}
\begin{itemize}
\raggedright
  \item \emph{Say it:} \emph{jhūṭhā}
  \item \emph{Say it:} \emph{jhūṭhā}, once more
  \item \emph{Say it:} say \emph{jhūṭhā}
  \item \emph{Recall:} say the Punjabi for raw, then the Punjabi for true, then say \emph{jhūṭhā} again
\end{itemize}

\begin{tcolorbox}[breakable,colback=teal!4,colframe=teal!35!black,title={Before you move on}]
What does \pa{ਝੂਠਾ} mean? (\textquotedblleft{}False, lying\textquotedblright{}.) Say it once more.
\end{tcolorbox}
